\bilingualchapter{Resources}{资源}

\bilingualsection{Further reading}{延伸阅读}

\bilingualsubsection{The flawed human brain}{有缺陷的人类大脑}

\bilingualsubsection{Thinking, Fast and Slow, Daniel Kahneman, 2011, Penguin}{《Thinking, Fast and Slow》，Daniel Kahneman，2011，Penguin}

\begin{englishblock} Excellent book on cognitive bias. A must read. \end{englishblock}

\begin{chineseblock}这是一本关于认知偏差的优秀著作，必读。\end{chineseblock}

\bilingualsubsection{Beyond Greed and Fear, Hersh Shefrin, 2007, OUP USA}{《Beyond Greed and Fear》，Hersh Shefrin，2007，OUP USA}

\begin{englishblock} Relatively short book on cognitive bias in finance specifically. Worth reading if you don't have time for Thinking, Fast and Slow. \end{englishblock}

\begin{chineseblock}这本书相对较短，专门讨论金融中的认知偏差。如果你没时间读 Thinking, Fast and Slow，它也非常值得一读。\end{chineseblock}

\bilingualsubsection{The Education of a Speculator, Victor Niederhoffer, 1998, Wiley}{《The Education of a Speculator》，Victor Niederhoffer，1998，Wiley}

\begin{englishblock} Fascinating and esoteric book by a famous partly systematic, negative skew, hedge fund manager. Part autobiography, part book on the philosophy of trading. Optional reading. \end{englishblock}

\begin{chineseblock}这本书出自一位著名的、部分系统化且带有负偏度风格的对冲基金经理之手，内容迷人而又带点冷门气质，一半像自传，一半像交易哲学书。属于可选阅读。\end{chineseblock}

\bilingualsubsection{Systematic trading rules}{系统化交易规则}

\bilingualsubsection{More Money than God, Sebastian Mallaby, 2011, Penguin}{《More Money than God》，Sebastian Mallaby，2011，Penguin}

\begin{englishblock} A history of hedge funds, but also a very readable guide to different strategies. Compulsory reading. \end{englishblock}

\begin{chineseblock}这是一本对冲基金史，同时也是一本非常易读的不同策略指南。必读。\end{chineseblock}

\bilingualsubsection{Expected Returns, Antti Ilmanen, 2011, Wiley Finance}{《Expected Returns》，Antti Ilmanen，2011，Wiley Finance}

\begin{englishblock} Comprehensive guide to the sources of returns and risk premia. Read after Mallaby, if you can cope with the sometimes technical treatment. \end{englishblock}

\begin{chineseblock}这是一本关于收益来源与风险溢价的全面指南。如果你能够接受其中有时较技术化的处理方式，建议在读完 Mallaby 之后再读这本。\end{chineseblock}

\bilingualsubsection{When Genius Failed: The Rise and Fall of Long-Term Capital Management, Roger Lowenstein, 2001, Random House}{《When Genius Failed: The Rise and Fall of Long-Term Capital Management》，Roger Lowenstein，2001，Random House}

\begin{englishblock} A negative skew disaster and a real-life example of the disadvantage of leveraging up too much on low volatility positions. Probably the most useful market history in this list and should be compulsory reading for anyone contemplating negative skew trading. \end{englishblock}

\begin{chineseblock}这是一个关于负偏度灾难的真实案例，也展示了在低波动头寸上加过多杠杆会带来怎样的坏处。它很可能是这份书单里最有价值的市场历史读物，任何考虑从事负偏度交易的人都应当把它列为必读。\end{chineseblock}

\bilingualsubsection{Rogue Trader, Nick Leeson, 1999, Sphere}{《Rogue Trader》，Nick Leeson，1999，Sphere}

\begin{englishblock} Very famous example of when negative skew goes wrong; amongst other things Nick lost much of his money selling option straddles. Optional reading. \end{englishblock}

\begin{chineseblock}这是一个非常著名的负偏度交易失控案例；Nick 在卖出跨式期权等操作中亏掉了大量资金。可选阅读。\end{chineseblock}

\bilingualsubsection{The Greatest Trade Ever, Gregory Zuckerman, 2010, Penguin}{《The Greatest Trade Ever》，Gregory Zuckerman，2010，Penguin}

\begin{englishblock} A great story about a positive skew trade that worked: John Paulson's bearish bet on mortgage backed securities. Optional reading. \end{englishblock}

\begin{chineseblock}这是一个关于成功正偏度交易的精彩故事：John Paulson 对抵押贷款支持证券所做的看空押注。可选阅读。\end{chineseblock}

\bilingualsubsection{The Black Swan, Nassim Taleb, 2008, Penguin}{《The Black Swan》，Nassim Taleb，2008，Penguin}

\begin{englishblock} A book about unknown unknowns. Taleb's usual mixture of unique philosophy and market folklore. Interesting, but optional reading. \end{englishblock}

\begin{chineseblock}这是一本讨论“未知的未知”的书。Taleb 一贯地把独特哲学与市场轶闻混合在一起。很有意思，但属于可选阅读。\end{chineseblock}

\bilingualsubsection{Trading rule fitting}{交易规则拟合}

\bilingualsubsection{Fooled by Randomness, Nassim Taleb, 2001, Penguin}{《Fooled by Randomness》，Nassim Taleb，2001，Penguin}

\begin{englishblock} A very interesting book on uncertainty in general. Compulsory for anyone who thinks back-testing is worthwhile. \end{englishblock}

\begin{chineseblock}这是一本关于一般性不确定性的非常有趣的著作。对于任何认为回测有价值的人来说，都是必读。\end{chineseblock}

\bilingualsubsection{Trading rules and forecasts}{交易规则与预测}

\bilingualsubsection{Trading Systems and Methods, 5th Edition, Perry J. Kaufman, 2013, John Wiley \& Sons}{《Trading Systems and Methods》第五版，Perry J. Kaufman，2013，John Wiley \& Sons}

\begin{englishblock} The bible of trading strategies. Great resource for trading rule ideas if you need them. \end{englishblock}

\begin{chineseblock}这是交易策略领域的“圣经”。如果你需要交易规则创意，它会是极好的资源。\end{chineseblock}

\bilingualsubsection{Technical Analysis, Jack Schwager, 1995, John Wiley \& Sons}{《Technical Analysis》，Jack Schwager，1995，John Wiley \& Sons}

\begin{englishblock} This, and the next book, are the best in a large crop of similar books. Compulsory if you want ideas for new technical trading rules. \end{englishblock}

\begin{chineseblock}这本书和下一本书，是同类大批作品中最好的两本。如果你想寻找新的技术交易规则创意，它们是必读。\end{chineseblock}

\bilingualsubsection{Fundamental Analysis, Jack Schwager, 1997, John Wiley \& Sons}{《Fundamental Analysis》，Jack Schwager，1997，John Wiley \& Sons}

\begin{englishblock} See above. Compulsory if you want to trade fundamentals. \end{englishblock}

\begin{chineseblock}同上。如果你打算交易基本面，它是必读。\end{chineseblock}

\bilingualsubsection{Hedge Fund Market Wizards, Jack Schwager, 2012, John Wiley \& Sons}{《Hedge Fund Market Wizards》，Jack Schwager，2012，John Wiley \& Sons}

\begin{englishblock} Interviews with many successful hedge fund managers. There are many useful nuggets of information in here. The chapter on Michael Platt is the most relevant to systematic traders. Optional reading. \end{englishblock}

\begin{chineseblock}本书采访了许多成功的对冲基金经理，其中有不少非常有用的信息碎片。关于 Michael Platt 的那一章，与系统化交易者最为相关。可选阅读。\end{chineseblock}

\bilingualsubsection{Volatility targeting}{波动率目标}

\bilingualsubsection{Fortune's Formula, William Poundstone, 2005, Hill \& Wang}{《Fortune's Formula》，William Poundstone，2005，Hill \& Wang}

\begin{englishblock} Highly readable story about the Kelly criterion, and the history of finance in general. Compulsory reading. \end{englishblock}

\begin{chineseblock}这是一本极具可读性的作品，讲述了 Kelly 准则以及更广义的金融史。必读。\end{chineseblock}

\audiencebox{Semi-automatic trader}{\textit{How to Win at Financial Spread Betting}, Charles Vincent, 2002, FT Prentice Hall. \par This is one of several books which provide a good introduction to the mechanics of spread betting in its first half. I don't advise using the strategies in the second half, but then I would say that!
\par
\medskip
\textit{How to Win at Financial Spread Betting}，Charles Vincent，2002，FT Prentice Hall。 \par 这本书的前半部分，是几本能很好介绍点差投注操作机制的书之一。至于后半部分的策略，我并不建议照着用，不过我想我本来也会这么说。}

\audiencebox{Asset allocating investor}{\textit{FT Guide to Exchange Traded Funds and Index Funds}, David Stevenson, 2012, Financial Times. \par UK book about ETFs.
\par
\medskip
\textit{The ETF Book}, Richard Ferri, 2009, Wiley. \par US book about ETFs.
\par
\medskip
\textit{FT Guide to Exchange Traded Funds and Index Funds}，David Stevenson，2012，Financial Times。 \par 这是一本文介绍 ETF 的英国书。
\par
\medskip
\textit{The ETF Book}，Richard Ferri，2009，Wiley。 \par 这是一本文介绍 ETF 的美国书。}

\audiencebox{Staunch systems trader}{\textit{Trading Commodities and Financial Futures}, George Kleinman, 2013, Prentice Hall. \par One of many books that give a good introduction to futures trading. There are also some ideas for trading rules. A fine alternative would be any of the Schwager books mentioned above.
\par
\medskip
\textit{Trading Commodities and Financial Futures}，George Kleinman，2013，Prentice Hall。 \par 这是众多能很好介绍期货交易入门的书之一，里面也包含一些交易规则思路。上面提到的任何一本 Schwager 著作，也都可以作为不错的替代选择。}

\bilingualsection{Sources of free data}{免费数据来源}

\begin{englishblock} For professional investors sourcing data is usually straightforward, but it can be harder for amateurs, especially those not wishing to pay. Sources of historical data that I've used are listed below. Website links are correct at the time of writing, but are subject to change. Many brokers and exchanges provide access to live prices, usually with a 15 minute delay. This delay is not problematic unless you are trading on a high frequency basis or using execution algorithms. \end{englishblock}

\begin{chineseblock}对于专业投资者来说，获取数据通常比较直接；但对于普通投资者，尤其是不想付费的人，这件事会更困难。下面列出了我自己使用过的一些历史数据来源。这些网站链接在写作时是正确的，但之后可能会发生变化。许多经纪商和交易所也提供实时价格访问，不过通常会有 15 分钟延迟。除非你在做高频交易或者使用执行算法，否则这种延迟并不会构成问题。\end{chineseblock}

\url{www.quandl.com}

\begin{englishblock} Source of numerous kinds of price and fundamental data, both free and subscription. Has a powerful API to automate data collection, available for multiple languages. For my own trading I currently use a combination of data from my broker and Quandl. \end{englishblock}

\begin{chineseblock}提供多种价格与基本面数据，既有免费也有订阅服务。它拥有强大的 API，可用于自动化数据采集，并支持多种语言。就我自己的交易而言，我目前使用的是经纪商数据与 quandl 数据的组合。\end{chineseblock}

\url{finance.yahoo.com}

\begin{englishblock} Source of equity and index data. Has an API to automate data collection. \end{englishblock}

\begin{chineseblock}提供股票和指数数据，也有可用于自动化数据采集的 API。 \end{chineseblock}

\url{www.bloomberg.com/markets/chart/data/1D/AAPL:US}

\begin{englishblock} One of the main providers of costly data for professional investors. As the link shows you can get a certain amount of free historic data from Bloomberg, if you know the reference code. \end{englishblock}

\begin{chineseblock} Bloomberg 是专业投资者昂贵数据服务的主要提供商之一。正如这个链接所展示的那样，如果你知道相应的参考代码，就可以从 Bloomberg 获取一定数量的免费历史数据。\end{chineseblock}

\url{www.eoddata.com}

\begin{englishblock} Free/subscription source of equity data. \end{englishblock}

\begin{chineseblock}提供免费和订阅制的股票数据。\end{chineseblock}

\url{www.ivolatility.com}

\begin{englishblock} Source of option prices. \end{englishblock}

\begin{chineseblock}提供期权价格数据。\end{chineseblock}

\url{www.oanda.com/currency/historical-rates}

\begin{englishblock} Source of historic FX rates. \end{englishblock}

\begin{chineseblock}提供历史外汇汇率数据。\end{chineseblock}

\url{stats.oecd.org}

\begin{englishblock} Official source of macroeconomic data. \end{englishblock}

\begin{chineseblock}官方宏观经济数据来源。\end{chineseblock}

\url{mba.tuck.dartmouth.edu/pages/faculty/ken.french/data_library.html}

\begin{englishblock} Academic source of equity value data. \end{englishblock}

\begin{chineseblock}股票价值数据的学术来源。\end{chineseblock}

\bilingualsection{Brokers and platforms}{经纪商与平台}

\begin{englishblock} If you are not running a fully automated trading system, I hesitate to recommend a specific broker. A key selection criteria is price. Estimate the expected annual total fee from how often you're trading, in what size, and with what total account value. Then use this for comparison purposes. Equally important is whether you can trade the products you want and with leverage if needed. Customer service should be a factor, but it is hard to evaluate in advance. Generally all other features are of secondary importance; although brokers seem to love jazzing up their platforms a flash website won't make you any more money. For those running an automated system there is only one broker accessible to amateur clients that I can recommend at the time of writing and that is Interactive Brokers (which I use myself ). This is because it offers a flexible API which can be used to fully automate \end{englishblock}

\begin{chineseblock}如果你运行的不是全自动交易系统，我不太愿意推荐某一家特定经纪商。选择时的一个关键标准是价格。你应根据自己的交易频率、交易规模以及总账户价值，估算预期的年度总费用，再据此进行比较。另一个同样重要的问题是：你能否交易自己想交易的产品，以及在需要时能否使用杠杆。客户服务也应纳入考虑，但这件事在事前很难真正评估。一般来说，其他功能都只是次要的；尽管经纪商似乎很喜欢把自己的平台包装得花哨漂亮，但一个炫目的网页并不会让你赚更多钱。对于运行自动化系统的人来说，在我写作本书时，我能推荐给普通客户的经纪商只有一家，那就是 Interactive Brokers（我自己也在用）。这是因为它提供了一个灵活的 API，可以用于完全自动化\end{chineseblock}

\bipairgap

\begin{englishblock} your trading and although it is not easy to get it working there is some help on my website: www.systematictrading.org There are online platforms that offer to implement your automated strategies for you, as distinct from software that runs on your desktop, which I discuss below. I haven't tested any of these but they look intriguing. Just be very careful, as you are handing over your money to someone else's software! There are also now ‘social trading' online platforms that allow you to put your money into other amateurs' systematic trading strategies. I would strongly advise against this. \end{englishblock}

\begin{chineseblock}你的交易。虽然把它真正跑起来并不轻松，但我的网站 \url{www.systematictrading.org} 上提供了一些帮助。市面上还有一些在线平台，可以替你部署自动化策略，这和运行在你自己桌面上的软件不同，后者我会在下面讨论。我没有亲自测试过这些平台，但它们看起来确实有点意思。不过一定要非常小心，因为你实际上是在把自己的钱交给别人的软件！现在还有一些“社交交易”在线平台，允许你把资金投入到其他普通人的系统化交易策略中。我强烈不建议这么做。\end{chineseblock}

\bilingualsection{Automation and coding}{自动化与编程}

\begin{englishblock} It is perfectly possible to run non-automated strategies using only simple spreadsheets. There are three avenues for implementing fully automated trading strategies. Firstly it is possible to use Excel to perform automated trading with Interactive Brokers, though this requires some knowledge of Visual Basic. Secondly you can use commercially available software, of which probably the most ubiquitous are NinjaTrader, TradeStation and MetaTrader, as well as broker specific packages. I haven't used, and so can't recommend, any particular product. The obvious advantage of these is that they allow you to create an automated strategy without having to write software. When deciding which to use, the most important criterion is the ability to implement your trading rules within the right framework. Back-testing technology is an added bonus, but make sure it is expanding out of sample or equivalently walk forward. Also it should allow you to fit across multiple instruments and include conservative cost estimates. Finally you can code up your own trading strategy. In the short term this is far more work, although it's also the most flexible and powerful option. The Interactive Brokers API allows you to use several languages, and there are third-party wrappers available which widen this choice further. For example I use Python to run my strategy via a library which is wrapped around the broker's C++ API. I include some snippets of python code on my website to show you how to implement various parts of the framework. \end{englishblock}

\begin{chineseblock}完全可以只借助简单的电子表格来运行非自动化策略。要实现全自动交易策略，大体有三条路径。第一，你可以使用 Excel 配合 Interactive Brokers 来进行自动化交易，不过这需要你掌握一些 Visual Basic 知识。第二，你可以使用市售软件，其中最常见的大概是 NinjaTrader、TradeStation 和 MetaTrader，以及各类经纪商自带的软件包。我自己没有使用过这些产品，因此无法推荐某一款具体工具。它们最明显的优势在于：你无需自己写软件，就可以创建自动化策略。在选择这类工具时，最重要的标准，是它能否让你在正确的框架中实现自己的交易规则。回测功能是额外加分项，但你必须确保它支持扩展样本外测试，或者等价的 walk-forward 方法。它还应当允许你在多个品种上进行拟合，并纳入保守的成本估计。最后，你也可以自己编写交易策略。从短期看，这会带来更多工作量，但它也是最灵活、最强大的选项。Interactive Brokers 的 API 允许你使用多种语言，而且还有第三方封装库可进一步扩大语言选择。比如，我自己就使用 Python，通过一个封装在该经纪商 C++ API 之上的库来运行策略。我在自己的网站上也放了一些 Python 代码片段，用来展示如何实现框架中的不同部分。\end{chineseblock}
